% =====================================================================
%  《药物使用者圣经》 LaTeX 排版模板 (preamble)
%  由 build_book.py 生成的 book.tex 引入。
%  可以自由修改本文件来调整版式；改完只需重新编译，无需重新生成正文。
%  编译引擎：XeLaTeX
% =====================================================================

% ---------- 页面 ----------
\usepackage{geometry}
\geometry{
  a4paper,
  top=2.5cm, bottom=2.6cm, left=2.7cm, right=2.7cm,
  headheight=15pt, headsep=14pt, footskip=1.1cm,
}

% ---------- 配色 ----------
\usepackage{xcolor}
\definecolor{dubaccent}{HTML}{2E6B78}    % 主色：墨青
\definecolor{dubaccentdk}{HTML}{1D454D}
\definecolor{dubtint}{HTML}{EDF4F5}      % 浅底
\definecolor{dubtint2}{HTML}{F7FAFA}
\definecolor{dubrule}{HTML}{B9CCD0}
\definecolor{dubink}{HTML}{23272A}
\definecolor{dubmute}{HTML}{6B7478}

% ---------- 字体 ----------
% 中文由 ctexbook 的 fontset=windows 提供：\songti \heiti \kaishu \fangsong
\usepackage{fontspec}
\IfFontExistsTF{TeX Gyre Termes}{%
  \setmainfont{TeX Gyre Termes}%
}{}
\IfFontExistsTF{TeX Gyre Heros}{%
  \setsansfont{TeX Gyre Heros}[Scale=MatchLowercase]%
}{}
\IfFontExistsTF{Consolas}{\setmonofont{Consolas}[Scale=MatchLowercase]}{}

\usepackage{microtype}

% ---------- 基础宏包 ----------
\usepackage{graphicx}
\usepackage{array}
\usepackage{booktabs}
\usepackage{longtable}
\usepackage{enumitem}
\usepackage{needspace}
\usepackage{ragged2e}
\usepackage{tabularray}
\UseTblrLibrary{booktabs}
% 正文里的表格没有题注，去掉 tabularray 自动生成的「表 N:」，并汉化跨页提示
\SetTblrTemplate{caption}{empty}
\SetTblrTemplate{caption-tag}{empty}
\SetTblrTemplate{caption-sep}{empty}
\SetTblrTemplate{caption-text}{empty}
\AtBeginDocument{%
  \DefTblrTemplate{contfoot-text}{normal}{\footnotesize\color{dubmute}（接下页）}%
  \DefTblrTemplate{conthead-text}{normal}{\footnotesize\color{dubmute}（接上页）}%
  \SetTblrTemplate{contfoot-text}{normal}%
  \SetTblrTemplate{conthead-text}{normal}%
}
\usepackage[most]{tcolorbox}
\usepackage{fancyhdr}
\usepackage{emptypage}

% 图片搜索路径（build_book.py 会把用到的图片链接/复制到 images/）
\graphicspath{{images/}}
\DeclareGraphicsExtensions{.jpeg,.jpg,.png,.pdf}

% ---------- 段落与断行 ----------
\setlength{\parskip}{0.42em plus 0.16em minus 0.08em}
\linespread{1.16}
\setlength{\emergencystretch}{3em}
\tolerance=2000
\hbadness=6000
\clubpenalty=8000
\widowpenalty=8000
\raggedbottom

% ---------- 章节标题样式 ----------
\ctexset{
  chapter = {
    name        = {第,章},
    number      = \arabic{chapter},
    format      = \raggedright,
    nameformat  = \normalsize\sffamily\color{dubaccent},
    aftername   = {\par\vskip6pt},
    titleformat = \huge\heiti\color{dubaccentdk},
    beforeskip  = 0pt,
    afterskip   = 28pt,
    aftertitle  = {\par\vskip8pt\noindent\textcolor{dubaccent}{\rule{\linewidth}{1.4pt}}\par},
    pagestyle   = plain,
  },
  section = {
    format     = \Large\heiti\color{dubaccentdk}\raggedright,
    aftername  = \enspace,
    beforeskip = 2.4ex plus .6ex minus .2ex,
    afterskip  = 1.2ex plus .2ex,
    aftertitle = {\par\vskip3pt\noindent\textcolor{dubrule}{\rule{\linewidth}{0.7pt}}\par},
  },
  subsection = {
    format     = \large\heiti\color{dubaccent}\raggedright,
    aftername  = \enspace,
    beforeskip = 2.6ex plus .6ex minus .2ex,
    afterskip  = 1.0ex plus .2ex,
  },
}

% ---------- 文件内小标题（不编号，但进入 PDF 书签） ----------
\makeatletter
\newcommand{\dubheadA}[1]{%
  \par\needspace{4\baselineskip}%
  \vspace{1.35em}%
  {\noindent\heiti\color{dubaccentdk}\large #1}%
  \par\vspace{0.18em}%
  {\noindent\textcolor{dubrule}{\rule{0.26\linewidth}{1.1pt}}}%
  \par\vspace{0.45em}%
  \phantomsection\addcontentsline{toc}{subsubsection}{#1}%
  \@afterindentfalse\@afterheading}
\newcommand{\dubheadB}[1]{%
  \par\needspace{3\baselineskip}%
  \vspace{1.0em}%
  {\noindent\heiti\color{dubaccent}\normalsize #1}%
  \par\vspace{0.3em}%
  \phantomsection\addcontentsline{toc}{paragraph}{#1}%
  \@afterindentfalse\@afterheading}
\newcommand{\dubheadC}[1]{%
  \par\needspace{2\baselineskip}%
  \vspace{0.7em}%
  {\noindent\heiti\color{dubink}\normalsize #1}%
  \par\vspace{0.22em}%
  \@afterindentfalse\@afterheading}
\makeatother

% ---------- 页眉页脚 ----------
\pagestyle{fancy}
\fancyhf{}
\renewcommand{\headrulewidth}{0.6pt}
\renewcommand{\footrulewidth}{0pt}
\renewcommand{\headrule}{\hbox to\headwidth{\color{dubrule}\leaders\hrule height \headrulewidth\hfill}}
\fancyhead[L]{\footnotesize\sffamily\color{dubmute}\nouppercase{\leftmark}}
\fancyhead[R]{\footnotesize\sffamily\color{dubmute}\nouppercase{\rightmark}}
\fancyfoot[C]{\footnotesize\sffamily\color{dubmute}\thepage}
\fancypagestyle{plain}{%
  \fancyhf{}%
  \renewcommand{\headrulewidth}{0pt}%
  \fancyfoot[C]{\footnotesize\sffamily\color{dubmute}\thepage}}
\renewcommand{\chaptermark}[1]{\markboth{\CTEXifname{\CTEXthechapter\quad}{}#1}{}}
\renewcommand{\sectionmark}[1]{\markright{\thesection\quad #1}}

% ---------- 插图 ----------
\newsavebox{\dubimgbox}
% 独立插图：保留原始尺寸，但不超过 0.70\linewidth / 0.40\textheight
\newcommand{\dubgraphic}[1]{%
  \sbox{\dubimgbox}{\includegraphics{#1}}%
  \ifdim\wd\dubimgbox>0.70\linewidth
    \includegraphics[width=0.70\linewidth,height=0.40\textheight,keepaspectratio]{#1}%
  \else
    \ifdim\ht\dubimgbox>0.40\textheight
      \includegraphics[height=0.40\textheight]{#1}%
    \else
      \usebox{\dubimgbox}%
    \fi
  \fi}
% 单张居中插图
\newcommand{\dubfig}[1]{%
  \par\vspace{0.7em}%
  \begin{center}\dubgraphic{#1}\end{center}%
  \vspace{0.25em}\par}
% 带说明文字的插图（长说明：两端对齐）
\newcommand{\dubfigcap}[2]{%
  \par\vspace{0.7em}%
  \begin{center}%
    \dubgraphic{#1}%
    \par\vspace{0.4em}%
    \begin{minipage}{0.88\linewidth}\small\color{dubmute}\setlength{\parskip}{0pt}#2\end{minipage}%
  \end{center}%
  \vspace{0.25em}\par}
% 带说明文字的插图（短说明：居中）
\newcommand{\dubfigcapc}[2]{%
  \par\vspace{0.7em}%
  \begin{center}%
    \dubgraphic{#1}%
    \par\vspace{0.4em}%
    \begin{minipage}{0.88\linewidth}\centering\small\color{dubmute}\setlength{\parskip}{0pt}#2\end{minipage}%
  \end{center}%
  \vspace{0.25em}\par}
% 并排插图：\dubrow{\dubinline{0.3}{a}\dubinline{0.3}{b}}
\newcommand{\dubrow}[1]{\par\vspace{0.7em}\noindent\begin{center}#1\end{center}\vspace{0.25em}\par}
\newcommand{\dubinline}[2]{%
  \begin{minipage}[c]{#1\linewidth}\centering
    \includegraphics[width=\linewidth,height=0.30\textheight,keepaspectratio]{#2}%
  \end{minipage}\hfill}
% 表格单元格内的插图
\newcommand{\dubcellimg}[1]{%
  \sbox{\dubimgbox}{\includegraphics{#1}}%
  \ifdim\wd\dubimgbox>\linewidth
    \includegraphics[width=\linewidth,height=6.5cm,keepaspectratio]{#1}%
  \else
    \ifdim\ht\dubimgbox>6.5cm
      \includegraphics[height=6.5cm]{#1}%
    \else
      \usebox{\dubimgbox}%
    \fi
  \fi}
% 源库中缺失的图片
\newcommand{\dubmissing}[1]{%
  \begin{center}\footnotesize\sffamily\color{dubmute}%
  \setlength{\fboxsep}{0pt}%
  \fcolorbox{dubrule}{dubtint2}{\parbox{0.55\linewidth}{\centering\vspace{9pt}〔图片缺失：\texttt{#1}〕\vspace{9pt}}}%
  \end{center}}

% ---------- 信息卡（药物速览表） ----------
\newenvironment{dubinfo}{%
  \par\vspace{0.6em}%
  \begin{tcolorbox}[
    enhanced, breakable,
    colback=dubtint2, colframe=dubtint2,
    boxrule=0pt, arc=1.4pt, outer arc=1.4pt,
    left=11pt, right=11pt, top=9pt, bottom=9pt, boxsep=0pt,
    borderline west={2.6pt}{0pt}{dubaccent},
  ]%
  \small
  \begin{tblr}{
    colspec={Q[l,mode=text,cmd=\heiti] X[l]},
    rowsep=2.8pt, colsep=7pt,
    hline{2-Y}={0.4pt, dubrule},
  }%
}{%
  \end{tblr}%
  \end{tcolorbox}\par\vspace{0.5em}}

% ---------- 引文 ----------
\newenvironment{dubquote}{%
  \par\vspace{0.55em}%
  \begin{tcolorbox}[
    enhanced, breakable, blanker,
    left=14pt, right=6pt, top=3pt, bottom=3pt,
    borderline west={2.2pt}{0pt}{dubrule},
  ]%
  \small\kaishu\color{dubaccentdk}\ignorespaces
}{%
  \end{tcolorbox}\par\vspace{0.45em}}

% ---------- 缩进块（源文件中的四空格缩进段落） ----------
\newenvironment{dubindent}{%
  \par\vspace{0.4em}%
  \begin{tcolorbox}[enhanced, breakable, blanker,
    left=16pt, right=6pt, top=2pt, bottom=2pt]%
  \small\setlength{\parskip}{0.25em}\ignorespaces
}{%
  \end{tcolorbox}\par\vspace{0.35em}}

% ---------- 列表 ----------
\setlist[itemize]{leftmargin=1.7em, itemsep=0.12em, parsep=0.12em, topsep=0.35em,
                  label=\textcolor{dubaccent}{\textbullet}}
\setlist[enumerate]{leftmargin=2.0em, itemsep=0.12em, parsep=0.12em, topsep=0.35em,
                    font=\color{dubaccent}\bfseries}

% ---------- 超链接 ----------
\usepackage[
  unicode,
  bookmarks=true, bookmarksnumbered=true, bookmarksopen=true, bookmarksopenlevel=1,
  bookmarksdepth=4,
  colorlinks=true,
  linkcolor=dubaccentdk, citecolor=dubaccentdk, urlcolor=dubaccent, filecolor=dubaccent,
  breaklinks=true, pdfstartview=FitH,
]{hyperref}
\usepackage{bookmark}
\usepackage{xurl}     % 长网址可以在任意位置断行，避免溢出页边
\urlstyle{same}

% 正文目录只到 subsection；PDF 书签更深
\setcounter{tocdepth}{2}
\setcounter{secnumdepth}{3}
